% =========================================================
% Proof Construction Algorithm
% =========================================================

\begin{remark}[What this section is for]
Architecture (Section~\ref{sec:proof-architecture}) tells you what kind
of proof to write. This section tells you how to write it, line by line,
once the architecture is chosen.

The procedure below is intentionally explicit and more detailed than
anything you will find in a standard textbook. That is on purpose.
With practice, these steps will become internalized and automatic.
But until they are --- until you have written enough proofs that the
process feels natural --- working through them consciously will prevent
the most common errors: undefined objects, unjustified steps, and
forgotten stop conditions.

A longer proof with every step explicit is always better than a shorter
proof with implicit gaps.
\end{remark}

\paragraph{Step 0: Classify the Statement.}

Before anything else, identify the logical form of the claim. Common
forms:
\begin{itemize}[itemsep=2pt]
  \item Universal: $\forall x\; P(x)$
  \item Conditional: $P \rightarrow Q$
  \item Biconditional: $P \leftrightarrow Q$
  \item Existential: $\exists x\; P(x)$
  \item Existence-and-uniqueness: $\exists!\, x\; P(x)$
  \item Set equality: $A = B$
  \item Algebraic equality: $a = b$
  \item Structural assertion: ``$f$ is injective,'' ``$R$ is an equivalence relation''
\end{itemize}

The logical form determines the skeleton of the proof. Consulting the
statement map (Section~\ref{sec:proof-architecture}) is the direct link
from this classification to an opening strategy.

\paragraph{Step 1: Restate the Givens.}

Write out explicitly what is given. Introduce every set, relation,
function, and algebraic structure, and record any properties assumed
about them.

\begin{quote}
Let $G$ be a group. Let $a \in G$.
\end{quote}

This is not a formality. By recording the givens, you license their
use in the proof --- you can now invoke associativity, the existence of
inverses, and the identity element without reintroducing them each time.
What is not in the givens cannot be used.

\paragraph{Step 2: Identify Objects and Their Types.}

Before reasoning begins, verify the \emph{type} of each object you
will work with: is it an element, a set, a function, or a relation?
What is its ambient universe? What operations apply to it?

\[
x \in A,\quad f : A \to B,\quad R \subseteq A \times A.
\]

This step prevents the most common logical errors in undergraduate
proofs: writing $x \in f$ when you mean $x \in \text{dom}(f)$; writing
$A = x$ when you mean $x \in A$; confusing membership ($\in$) with
inclusion ($\subseteq$). Every object used in a proof must have a
declared type. If you cannot say what type something is, you do not
yet understand what you are working with.

\paragraph{Step 3: Introduce Arbitrary Elements.}

If the claim is universal ($\forall x\; P(x)$), immediately introduce
an arbitrary element:
\[
\text{Let } x \in A \text{ be arbitrary.}
\]

The word ``arbitrary'' is not decorative. It means the element is not
given any properties beyond those of the set $A$. Whatever you prove
about this $x$ will hold for all $x \in A$, precisely because you have
assumed nothing about it beyond membership. The moment you specialize
--- ``let $x = 0$'' or ``let $x$ be such that \ldots'' --- you have
lost the ability to conclude the universal statement.

\paragraph{Step 4: Expand the Goal.}

Rewrite the conclusion using definitions rather than named concepts.
\begin{itemize}[itemsep=2pt]
  \item To prove $f$ is injective: expand injectivity as
    ``$f(a) = f(b) \Rightarrow a = b$'' and take that as the new goal.
  \item To prove $A = B$ (sets): take as new goal
    ``$A \subseteq B$ and $B \subseteq A$.''
  \item To prove $x \in A \cup B$: rewrite the goal as ``$x \in A$
    or $x \in B$.''
  \item To prove $\{x_n\}$ converges: expand to
    ``$\forall\, \varepsilon > 0\; \exists\, N$ s.t.\ $n \geq N
    \Rightarrow d(x_n, L) < \varepsilon$.''
\end{itemize}

This step is the one most students skip, and skipping it is the single
most reliable predictor of a stuck proof. You cannot make progress toward
a goal you have not named in operative terms. ``Prove $f$ is injective''
is not an operative goal. ``Prove that $f(a) = f(b)$ implies $a = b$''
is.

\paragraph{Step 5: Introduce Helper Objects.}

Introduce any auxiliary objects needed for the argument: witnesses
for existential claims, intermediate elements for indirect arguments,
bounds for inequality proofs, function values.

\[
\text{Let } N_1 \in \mathbb{N} \text{ s.t.\ } n \geq N_1
\Rightarrow d(x_n, L) < \varepsilon/2.
\]

No object should appear in a proof without being explicitly introduced.
If you write ``let $c$ be the cancellation element'' without first
establishing that such an element exists, your proof has an undefined
object. Every name carries an obligation: the obligation to have said
what the name refers to.

\paragraph{Step 6: Apply One Rule at a Time.}

Each line of a proof follows from exactly one of four sources:
\begin{enumerate}[itemsep=2pt]
  \item A definition (unpacking what a term means),
  \item A hypothesis (something assumed in the current proof),
  \item A previously proved theorem (cited by name and reference),
  \item A basic logical rule (modus ponens, universal instantiation, etc.).
\end{enumerate}

Lines that are justified by ``clearly'' or ``obviously'' or no reason at
all are gaps. At the learning stage, every step should have an explicit
justification. When a proof stalls, the right question is always: which
of the four sources has not been consulted? Almost always, the answer is
that some definition has not been unpacked.

\paragraph{Step 7: Use Forward and Backward Reasoning.}

Proofs alternate between two modes of reasoning, and recognizing which
mode is productive at each stage is a significant skill.

\emph{Forward reasoning:} deduce consequences from what you have
established so far. If you know $a^2 = e$ in a group, derive that
$a(ae) = a \cdot a \cdot e = a^2 \cdot e = e$. You are moving away
from the hypotheses, toward the goal.

\emph{Backward reasoning:} examine the goal and ask what would
suffice to establish it. If the goal is ``$a = a^{-1}$,'' ask: by
what definition would that follow? It follows if $aa = e$, since
that is the condition for $a$ to be its own inverse. Now your new
goal is to prove $aa = e$, which may be more directly accessible.

In practice, most proofs proceed in both directions simultaneously:
you derive forward from the hypotheses while simplifying the goal
backward until the two sides meet. When you are stuck, try the
direction you have not been working in.

\paragraph{Step 8: Handle Cases Explicitly.}

If a statement splits into mutually exclusive and exhaustive cases,
enumerate all of them.

\[
\text{Either } n \text{ is even, or } n \text{ is odd.}
\]

Handle each case separately. Close each case before opening the next.
Verify, explicitly, that the cases together cover every possibility.
A proof with unlabeled cases --- where it is unclear which case is
being argued or whether all cases have been covered --- is not a proof.

\paragraph{Step 9: Close the Argument.}

Once the desired conclusion has been established, state it explicitly:

\begin{itemize}[itemsep=2pt]
  \item ``Thus $x \in B$, and so $A \subseteq B$.''
  \item ``Hence $f(a) = f(b)$ implies $a = b$, so $f$ is injective.''
  \item ``In both cases $P(n+1)$ holds.''
\end{itemize}

The explicit closing step serves two functions. First, it tells the reader
(and you) that the proof goal has been achieved. Second, it forces you to
check that what you have proved actually \emph{is} the goal. Students who
reach the end of a proof and write \qed without stating a conclusion often
discover, on re-reading, that they proved something adjacent to the goal
but not the goal itself.

\paragraph{Step 10: Signal Completion.}

Conclude with \qed, $\blacksquare$, or ``This completes the proof.''

\bigskip

\begin{tcolorbox}[colback=propbox, colframe=propborder, arc=2pt,
  left=6pt, right=6pt, top=4pt, bottom=4pt,
  title={\small\textbf{Legal Moves and Stop Conditions}},
  fonttitle=\small\bfseries]
\textbf{What you may never do:}
\begin{itemize}[itemsep=2pt]
  \item Choose an existential witness before proving it exists.
  \item Introduce an ``arbitrary'' element and then give it special properties.
  \item Assume the conclusion at any point in the proof.
  \item Apply a definition partially (all conditions must be checked).
\end{itemize}

\medskip
\textbf{A proof is complete when:}
\begin{itemize}[itemsep=2pt]
  \item A universal claim has been proved for a genuinely arbitrary element.
  \item An existential claim has produced a verified witness.
  \item A set equality has established both inclusions explicitly.
  \item Each direction of a biconditional has been proved separately.
  \item An existence-uniqueness claim has done both parts.
\end{itemize}
\end{tcolorbox}

\paragraph{Line-by-Line Discipline.}

Before writing each line, ask three silent questions:
\begin{enumerate}[label=(\roman*), itemsep=2pt]
  \item Has every object mentioned in this line been defined or introduced earlier?
  \item Which definition, hypothesis, or theorem justifies this step?
  \item Does this step move the argument toward the stated goal?
\end{enumerate}
If any answer is no, stop and resolve it before continuing.
